\theory{Surgery theory}{How can one construct or classify manifolds by cutting out a controlled part and replacing it with another part?}{Surgery theory modifies a manifold along an embedded sphere, keeping track of the changes in homology, homotopy, and other invariants. In high dimensions it turns manifold classification into a sequence of obstruction problems.}{High-dimensional topology, cobordism, exotic spheres, manifold classification, and geometric topology.}{Cutting out a tubular neighborhood and regluing it changes a manifold in a controlled way whose homotopy and obstruction data can be computed.}

\paragraph{The operation and history}
For an $n$-manifold, remove an embedded $S^p$ with a tubular neighborhood $S^p\times D^q$ and glue in $D^{p+1}\times S^{q-1}$, where $n=p+q+1$. The boundary is the same, so the replacement produces another manifold. John Milnor called the method surgery; Andrew Wallace called it spherical modification \cite{surgery_ref1}.

Surgery arose from the classification of exotic spheres and became a central tool for manifolds of dimension greater than four. Morse theory shows that two manifolds are related by a sequence of such modifications exactly when they lie in the same cobordism class \cite{surgery_ref2}.

\end{historylist}
\section*{3. Theory}
\subsection*{Core theory}
\paragraph{What surgery controls}
The embedded sphere determines which homotopy class is changed. A surgery can kill a homotopy group, alter homology in a known degree, or simplify the fundamental group. The goal is to start with a manifold whose structure is known and produce one satisfying a desired property.

The main classification questions are whether a space with the homotopy type of a manifold is actually homeomorphic or diffeomorphic to one, and whether a homotopy equivalence between manifolds is homotopic to a diffeomorphism. For dimensions above four, surgery gives primary and secondary obstructions rather than a single complete invariant \cite{surgery_ref3}.

\paragraph{Worked example: surgery on $S^3$}
Consider the three-sphere $S^3$ and perform surgery on an embedded circle $S^1\subset S^3$. The normal bundle of $S^1$ in $S^3$ is trivial (since $S^3$ is parallelizable), so a tubular neighborhood looks like $S^1\times D^2$. Surgery means removing this $S^1\times D^2$ and gluing in $D^2\times S^1$ along the common boundary $S^1\times S^1$. The result is again $S^3$: one can see this by observing that $S^3$ is the boundary of $D^4$, and the surgery corresponds to attaching a 2-handle along the unknot, which does not change the manifold. More interestingly, surgery on a knotted $S^1\subset S^3$ (a nontrivial knot) can produce exotic manifolds. For example, surgery on the right-hand trefoil knot with framing $1$ yields the Poincaré homology sphere, a manifold with the same homology as $S^3$ but with $\pi_1\cong$ the binary icosahedral group of order 120. This shows that even in dimension three, surgery can change the topology in ways invisible to homology.

\paragraph{Obstructions and the s-cobordism theorem}
A normal map records a candidate manifold map together with stable normal bundle data. A primary obstruction in topological $K$-theory must vanish before a surgery obstruction can be defined. The secondary obstruction lies in an algebraic $L$-group. When it vanishes, surgery can produce the desired homotopy equivalence under suitable hypotheses.

The h-cobordism theorem says that an h-cobordism between simply connected manifolds of dimension at least five is a product. The s-cobordism theorem removes simple homotopy ambiguity and uses Whitehead torsion. These results explain why high-dimensional classification is tractable while dimensions three and four require different methods \cite{surgery_ref4}.

\section*{4. Important theorems and results}
\begin{enumerate}[leftmargin=*,itemsep=0.35em]
\item \textbf{Surgery construction.} Replacing $S^p\times D^q$ by $D^{p+1}\times S^{q-1}$ produces a controlled change of an $n$-manifold.

\item \textbf{Cobordism-surgery principle.} A cobordism can be decomposed into elementary surgeries.

\item \textbf{h-cobordism theorem.} A simply connected h-cobordism in dimension at least five is a product.

\item \textbf{s-cobordism theorem.} In the high-dimensional setting, vanishing Whitehead torsion characterizes when an h-cobordism is a product.

\item \textbf{Surgery exact sequence.} Manifold structures, normal invariants, and algebraic $L$-groups are linked by an exact sequence of classification problems.
\end{enumerate}
\noindent\emph{Source for these results:} \cite{surgery_ref1}.

\theoryapplications
\theoryconnections{surgery_theory}{%
\item Surgery removes a tubular neighborhood of an embedded sphere in a manifold and replaces it with another piece having the same boundary.
\item Homotopy and homology calculate how this replacement changes the space, cobordism records the whole transformation, and vector bundles describe the normal data needed to perform it.
\item K-theory and L-theory turn the remaining geometric question into an algebraic obstruction, especially when a homotopy equivalence is to be improved to a genuine equivalence of manifolds \cite{surgery_ref1,surgery_ref2}.
}{%
\item Surgery theory classifies high-dimensional manifolds by reducing them to a sequence of controlled modifications.
\item It explains exotic spheres by showing when the standard sphere can be altered without changing its homotopy type, and it supplies the obstruction groups used in the h- and s-cobordism theorems.
\item The connection is powerful but dimension-sensitive: the Whitney trick and general-position arguments that make the algebra reliable are much weaker in low dimensions \cite{surgery_ref3,surgery_ref4}.
\item \textbf{Three-manifold topology and the Poincar\'e conjecture.} Surgery theory was the conceptual engine behind the resolution of the Poincar\'e conjecture. Perelman's proof via Ricci flow can be viewed as a continuous, analytic version of surgery: as the flow progresses, singularities develop that are modeled on neck-pinch operations, and each singularity is resolved by a surgery that removes and replaces a cylindrical region. The classification program for three-manifolds---the Thurston geometrization conjecture, now a theorem---decomposes every closed 3-manifold along incompressible spheres and tori into pieces that each carry one of eight model geometries. Surgery-theoretic language (cutting along embedded spheres, tracking the effect on $\pi_1$, and using obstruction theory) provides the combinatorial framework that organizes this decomposition and led to Perelman's proof.
}
\section*{Glossary}
\begin{description}[leftmargin=*,style=nextline,itemsep=0.3em]
\item[\textbf{Surgery.}] A modification of a manifold by removing a tubular neighborhood of an embedded sphere and gluing in a different piece with the same boundary.
\item[\textbf{Tubular neighborhood.}] A neighborhood of a submanifold that looks like a disk bundle over the submanifold; in surgery it is $S^p\times D^q$.
\item[\textbf{Cobordism.}] A relationship between manifolds where one is the boundary of a higher-dimensional manifold; surgery decomposes cobordisms into elementary steps.
\item[\textbf{h-cobordism theorem.}] States that an h-cobordism between simply connected manifolds of dimension at least five is a product.
\item[\textbf{s-cobordism theorem.}] Refines the h-cobordism theorem by using Whitehead torsion to characterize when an h-cobordism is a product.
\item[\textbf{Whitehead torsion.}] An algebraic invariant that measures the failure of a homotopy equivalence to be a simple homotopy equivalence.
\item[\textbf{Normal map.}] A map between manifolds together with stable normal bundle data used to define surgery obstructions.
\item[\textbf{L-group.}] An algebraic group that classifies the secondary obstructions in surgery theory.
\item[\textbf{Exotic sphere.}] A smooth manifold homeomorphic but not diffeomorphic to the standard sphere.
\item[\textbf{K-theory.}] An algebraic theory of stable equivalence classes of vector bundles providing the primary obstruction group.
\item[\textbf{Whitney trick.}] A topological technique for removing intersection points using an embedded disk, available in high dimensions.
\item[\textbf{Normal bundle.}] The vector bundle of directions perpendicular to an embedded submanifold.
\end{description}
\section*{7. Sample Questions}
\emph{Exercise reference:} \cite{surgery_ref1}. The cited edition is the source for the exercise set; consult its Exercises section for the official statement and numbering.
\begin{enumerate}[leftmargin=*,itemsep=0.9em]
\item \textbf{Exercise 1.} In surgery on a 4-manifold $M^4$, suppose we perform surgery on an embedded $S^1 \subset M^4$. Identify $p$, $q$, and the pieces removed and glued in. \emph{Solution.} For $n = 4 = p + q + 1$ with $p = 1$, we get $q = 2$. The tubular neighborhood removed is $S^1 \times D^2$, and the piece glued in is $D^2 \times S^1$. Both have boundary $S^1 \times S^1$, so the gluing is along a torus.

\item \textbf{Exercise 2.} Perform surgery on the unknot $S^1 \subset S^3$. Compute the Euler characteristic before and after surgery. \emph{Solution.}$S^3$ is a closed 3-manifold with $\chi(S^3) = 0$ (since $H_0 = H_3 = \mathbb{Z}$ and $H_1 = H_2 = 0$, by Poincaré duality). Surgery on the unknot in $S^3$ produces $S^2 \times S^1$ (removing $S^1 \times D^2$ and gluing $D^2 \times S^1$). We compute $\chi(S^2 \times S^1) = \chi(S^2) \cdot \chi(S^1) = 2 \cdot 0 = 0$. The Euler characteristic is unchanged, consistent with the general fact that surgery on an $S^p$ does not change $\chi$.

\item \textbf{Exercise 3.} Compute the Betti numbers of the Poincaré homology sphere $\Sigma(2,3,5)$, which can be obtained by $(+1)$-surgery on the right-hand trefoil knot in $S^3$. \emph{Solution.}$\Sigma(2,3,5)$ is a rational homology sphere: $H_0 = \mathbb{Z}$, $H_1 = 0$ (the first homology of a Seifert fibered space over $S^2$ with relatively prime invariants), $H_2 = 0$ (by Poincaré duality with integer coefficients), and $H_3 = \mathbb{Z}$. The Betti numbers are $b_0 = 1$, $b_1 = 0$, $b_2 = 0$, $b_3 = 1$. The fundamental group is the binary icosahedral group of order 120, confirming $H_1 = \pi_1/[\pi_1, \pi_1] = 0$.

\item \textbf{Exercise 4.} In a surgery on $S^3$ along a knotted $S^1$, the framing determines the resulting manifold. Explain why surgery on the unknot with framing $n$ produces $S^2 \times S^1$ for any integer $n$ (up to orientation). \emph{Solution.} The unknot bounds an embedded disk $D^2$ in $S^3$. Surgery on the unknot with framing $n$ corresponds to attaching a 2-handle to $D^4$ along the unknot with framing $n$. Since the unknot has trivial knot group, the handle attachment does not create any nontrivial topology: it simply thickens the disk. The result is always $S^2 \times S^1$ (up to orientation), independent of $n$, because the surgery on the unknot is equivalent to a connected sum decomposition that produces the same manifold.

\item \textbf{Exercise 5.} Verify that the h-cobordism theorem requires dimension $\geq 5$ by explaining why the Whitney trick fails in dimension 4. \emph{Solution.}$H_1(X; \mathbb{Z}) \cong \mathbb{Z}$ has rank 1 (non-zero first homology). In the construction of the Whitney disk in 4-manifolds, the disk's interior can self-intersect because $\pi_1(X)$ is non-trivial. Specifically, the two arcs forming the Whitney disk boundary determine an element in $\pi_1$, and in dimension 4 the disk cannot avoid self-intersecting when this element is non-trivial. This means the Whitney trick fails to remove pairs of intersection points in dimension 4, unlike in dimension $\geq 5$ where there is enough room to make the disk embedded. This is why the h-cobordism theorem and surgery classification require dimension $\geq 5$.

\end{enumerate}
\section*{8. References}
\begin{thebibliography}{9}
\bibitem{surgery_ref1} J. Milnor, \emph{Collected Papers, Volume III: Differential Topology}, American Mathematical Society, 2010.
\bibitem{surgery_ref2} W. Browder, \emph{Surgery on Simply-Connected Manifolds}, Springer-Verlag, 1972.
\bibitem{surgery_ref3} C. T. C. Wall, \emph{Surgery on Compact Manifolds}, American Mathematical Society, 1999.
\bibitem{surgery_ref4} A. Ranicki, \emph{Algebraic and Geometric Surgery}, Oxford University Press, 2002.
\end{thebibliography}
